% MT047 / explanation — canonical editable content.
题中每条记录有输入\(x_i\)和标签\(y_i\in\{-1,+1\}\)。模型先算\(f(x_i)=w^Tx_i+b\)；它可正可负，也可大于1，因此\textbf{不是概率}。把分数乘真实标签，得到\(z_i=y_if(x_i)\)：\(z_i>0\)说明方向正确，\(z_i<0\)说明方向错误。


用补充数字看Logistic损失\(\ell(z)=\log(1+e^{-z})\)怎样工作：正类\(y=1\)若得到\(f=2\)，\(z=2\)，损失约0.127；若得到\(f=-2\)，\(z=-2\)，损失约2.127。负类\(y=-1\)若得到\(f=-2\)，乘积同样为2，也受到小惩罚。因此同一个公式同时处理正负两类。


sigmoid给真实标签的概率为\(\sigma(z_i)=1/(1+e^{-z_i})\)。这里\(e\approx2.718\)，\(\log\)是以\(e\)为底的自然对数；\(\log(1/u)=-\log u\)，所以取负log得到



\[
-\log\frac{1}{1+e^{-z_i}}=\log(1+e^{-z_i}).
\]


概率越接近1，负log越接近0；概率越接近0，负log越大。因此这个损失会惩罚“给真实类别很低概率”的预测。


求和\(\sum_i\ell(z_i)\)就是把所有样本的代价加起来。减小它，意味着整体上给真实标签更高概率。


另一项\(\alpha w^Tw=\alpha\sum_jw_j^2\)对大权重收费。例如同样把损失降低一点，用巨大的系数通常比用较小系数付出更多正则代价。\(\alpha\)越大，模型越倾向缩小权重；太大可能连真实规律也学不出，产生欠拟合。\(\alpha\)很小则更着重贴合训练数据，可能过拟合，应用验证选择合适值。


本式\(\alpha\)乘在惩罚前，另一种课堂写法以\(1/C\)乘惩罚；在这些约定下方向相反。不要只背“参数变大就更复杂”，先看它在目标中乘哪一项。


选读：零均值高斯参数先验的负log为平方惩罚，故\(\alpha\)与先验方差的倒数成正比。按题目使用损失求和还是平均，具体比例常数也会变化。


\textbf{出处：} 2023A期中 Q8，原卷第5页（10分）；2025A期中 Q8，原卷第5页（10分）


\textbf{答案依据：} 2023A期中Q8随卷答案；2025A Q8同题。分数与概率的区别已复核。
